% Options for packages loaded elsewhere
\PassOptionsToPackage{unicode}{hyperref}
\PassOptionsToPackage{hyphens}{url}
\PassOptionsToPackage{dvipsnames,svgnames,x11names}{xcolor}
\documentclass[
  10pt,
]{article}
\usepackage{xcolor}
\usepackage[margin=2.4cm]{geometry}
\usepackage{amsmath,amssymb}
\setcounter{secnumdepth}{-\maxdimen} % remove section numbering
\usepackage{iftex}
\ifPDFTeX
  \usepackage[T1]{fontenc}
  \usepackage[utf8]{inputenc}
  \usepackage{textcomp} % provide euro and other symbols
\else % if luatex or xetex
  \usepackage{unicode-math} % this also loads fontspec
  \defaultfontfeatures{Scale=MatchLowercase}
  \defaultfontfeatures[\rmfamily]{Ligatures=TeX,Scale=1}
\fi
\usepackage{lmodern}
\ifPDFTeX\else
  % xetex/luatex font selection
  \setmainfont[]{Libertinus Serif}
  \setmonofont[Scale=0.85]{Latin Modern Mono}
  \setmathfont[]{Latin Modern Math}
\fi
% Use upquote if available, for straight quotes in verbatim environments
\IfFileExists{upquote.sty}{\usepackage{upquote}}{}
\IfFileExists{microtype.sty}{% use microtype if available
  \usepackage[]{microtype}
  \UseMicrotypeSet[protrusion]{basicmath} % disable protrusion for tt fonts
}{}
\makeatletter
\@ifundefined{KOMAClassName}{% if non-KOMA class
  \IfFileExists{parskip.sty}{%
    \usepackage{parskip}
  }{% else
    \setlength{\parindent}{0pt}
    \setlength{\parskip}{6pt plus 2pt minus 1pt}}
}{% if KOMA class
  \KOMAoptions{parskip=half}}
\makeatother
\usepackage{longtable,booktabs,array}
\newcounter{none} % for unnumbered tables
\usepackage{calc} % for calculating minipage widths
% Correct order of tables after \paragraph or \subparagraph
\usepackage{etoolbox}
\makeatletter
\patchcmd\longtable{\par}{\if@noskipsec\mbox{}\fi\par}{}{}
\makeatother
% Allow footnotes in longtable head/foot
\IfFileExists{footnotehyper.sty}{\usepackage{footnotehyper}}{\usepackage{footnote}}
\makesavenoteenv{longtable}
\setlength{\emergencystretch}{3em} % prevent overfull lines
\providecommand{\tightlist}{%
  \setlength{\itemsep}{0pt}\setlength{\parskip}{0pt}}

\providecommand{\E}{\mathbb{E}}
\providecommand{\N}{\mathbb{N}}
\providecommand{\Z}{\mathbb{Z}}
\providecommand{\R}{\mathbb{R}}
\providecommand{\eps}{\varepsilon}
\providecommand{\ceil}[1]{\left\lceil #1\right\rceil}
\providecommand{\floor}[1]{\left\lfloor #1\right\rfloor}
\AtBeginDocument{\let\setminus\smallsetminus}
\usepackage[most]{tcolorbox}
\newtcolorbox{warning}{colback=red!5,colframe=red!65!black,boxrule=0.8pt,arc=2pt,left=6pt,right=6pt,top=5pt,bottom=5pt}
\usepackage{etoolbox}
\AtBeginEnvironment{longtable}{\small}
\setlength{\emergencystretch}{3em}
\usepackage{fancyhdr}
\pagestyle{fancy}
\fancyhf{}
\fancyhead[L]{\small\bfseries UNREFEREED GPT-6 ASTRA OUTPUT}
\fancyhead[R]{\small\thepage}
\renewcommand{\headrulewidth}{0.4pt}
\usepackage{bookmark}
\IfFileExists{xurl.sty}{\usepackage{xurl}}{} % add URL line breaks if available
\urlstyle{same}
\hypersetup{
  pdftitle={Conjecture 24},
  pdfauthor={Writeup: gpt-6-astra (single model pass)},
  colorlinks=true,
  linkcolor={blue!50!black},
  filecolor={Maroon},
  citecolor={Blue},
  urlcolor={blue!50!black},
  pdfcreator={LaTeX via pandoc}}

\title{Conjecture 24}
\usepackage{etoolbox}
\makeatletter
\providecommand{\subtitle}[1]{% add subtitle to \maketitle
  \apptocmd{\@title}{\par {\large #1 \par}}{}{}
}
\makeatother
\subtitle{Unrefereed candidate disproof by GPT-6 Astra}
\author{Writeup: \texttt{gpt-6-astra} (single model pass)}
\date{Catalog id \texttt{2506.08810\_\_03} --- generated 2026-09-17}

\begin{document}
\maketitle

\begin{warning}

\textbf{UNREFEREED MODEL OUTPUT.} This document was selected solely
because GPT-6 Astra labelled its own result
\texttt{would\_publish:\ true} and returned \texttt{proved} or
\texttt{disproved}. It has not passed the adversarial LLM referee used
for the earlier Sol campaign, has not been checked by a human
mathematician, and has not been checked for novelty. Treat every
mathematical and bibliographic claim below as unverified.

\end{warning}

\section{Summary and provenance}\label{summary-and-provenance}

{\def\LTcaptype{none} % do not increment counter
\begin{longtable}[]{@{}
  >{\raggedright\arraybackslash}p{(\linewidth - 2\tabcolsep) * \real{0.5000}}
  >{\raggedright\arraybackslash}p{(\linewidth - 2\tabcolsep) * \real{0.5000}}@{}}
\toprule\noalign{}
\endhead
\bottomrule\noalign{}
\endlastfoot
Catalog id & \texttt{2506.08810\_\_03} \\
Catalog entry &
\href{https://graph-theory-ai.github.io/graph-conjectures/arxiv/2506.08810__03/}{Conjecture
24} \\
Source corpus & arxiv \\
Campaign leg & selected second attempts (\texttt{attacks\_retry}) \\
Source paper / entry & Infinite induced-saturated graphs \\
Model verdict & \textbf{disproved} (confidence: high) \\
Model's one-line claim & The five-vertex tournament C3{[}TT2,TT2,1{]} is
a counterexample: every countable tournament avoiding it has a distinct
locally finite perturbation that still avoids it. \\
Model & \texttt{gpt-6-astra}, reasoning effort \texttt{max},
\texttt{mode=pro}, flex service tier \\
Original artifact &
\texttt{attacks\_retry/2506.08810\_\_03/output.md} \\
Independent review & \textbf{None} \\
Model's caveats & Perturbations are understood to be nonempty; the
argument does not settle the weaker single-arc saturation question. \\
\end{longtable}
}

\section{Problem statement}\label{problem-statement}

\emph{No extracted statement was stored in the artifact metadata.}

\section{Astra writeup}\label{astra-writeup}

\emph{The text below is the model output verbatim, apart from moving its
machine-readable verdict block into the summary above and shifting
Markdown heading levels for this document. Mechanical TeX defects and
equation tags were normalized where needed for compilation.}

\subsection{1. The counterexample}\label{the-counterexample}

Let \(T_\ast\) be the five-vertex tournament with vertex partition \[
A=\{a_0,a_1\},\qquad B=\{b_0,b_1\},\qquad \{c\},
\] and orientations \[
a_0\to a_1,\qquad b_0\to b_1,\qquad
A\to B\to \{c\}\to A.
\] Thus \[
T_\ast=C_3[\mathrm{TT}_2,\mathrm{TT}_2,1].
\] Here, an arrow between sets means that all arcs between those sets
have the indicated direction.

This tournament is nontransitive. I will prove the following, which
directly contradicts the conjecture.

\subsubsection{Theorem}\label{theorem}

For every countably infinite \(T_\ast\)-free tournament \(S\), there is
a \(T_\ast\)-free tournament \(S'\ne S\) on the same vertex set such
that the pairs whose orientations differ between \(S\) and \(S'\) form a
locally finite graph.

The key structural fact is stronger than needed:

\begin{quote}
Every countable \(T_\ast\)-free tournament can be made transitive by a
locally finite set of arc reversals.
\end{quote}

No result from the previous attempt is needed below.

\subsection{2. A criterion for locally finite distance from a transitive
tournament}\label{a-criterion-for-locally-finite-distance-from-a-transitive-tournament}

For a tournament \(S\), define an undirected graph \(J(S)\) on \(V(S)\)
by declaring \(xy\) to be an edge precisely when the arc between \(x\)
and \(y\) belongs to infinitely many directed triangles.

Write \[
O(x)=N_S^+(x).
\] For an arc \(x\to y\), its directed-triangle completions are exactly
\[
Z(x,y)=\{z:y\to z\to x\}.
\] In particular, \[
O(y)\setminus O(x)=Z(x,y).
\qquad\text{(1)}
\]

\subsubsection{Lemma 1}\label{lemma-1}

A countable tournament \(S\) differs from some transitive tournament on
\(V(S)\) by a locally finite set of arc reversals if and only if
\(J(S)\) is locally finite.

\paragraph{Proof}\label{proof}

\textbf{Necessity.} Suppose \(L\) is transitive and the change graph
\(D=D(S,L)\) is locally finite.

If \(xy\notin E(D)\), then every directed triangle \(xyz\) of \(S\) must
have \(xz\in E(D)\) or \(yz\in E(D)\): otherwise that triangle would
also occur in \(L\). Thus there are only finitely many directed
triangles of \(S\) containing \(xy\). Consequently, \[
J(S)\subseteq D,
\] so \(J(S)\) is locally finite.

\textbf{Sufficiency.} Suppose \(J(S)\) is locally finite. For sets
\(X,Y\), write \[
X\subseteq^\ast Y
\quad\Longleftrightarrow\quad
X\setminus Y\text{ is finite}.
\] Define an equivalence relation on \(V(S)\) by \[
x\sim y
\quad\Longleftrightarrow\quad
O(x)\mathbin{\triangle}O(y)\text{ is finite}.
\]

On the equivalence classes, reverse almost-inclusion of
outneighbourhoods defines a partial order: a class \(C\) precedes a
distinct class \(D\) when \[
O(y)\subseteq^\ast O(x)
\qquad(x\in C,\ y\in D).
\] This is well-defined, and antisymmetry follows from the definition of
\(\sim\). Choose a linear extension of this partial order.

If \(x\to y\) and \(xy\notin E(J(S))\), then (1) gives \[
O(y)\subseteq^\ast O(x).
\] Therefore, if \(x\) and \(y\) belong to different equivalence
classes, their original orientation agrees with the chosen class order.
Hence:

\[
\text{Every incorrectly oriented pair between classes belongs to }J(S).
\qquad\text{(2)}
\]

It remains to order each equivalence class with only locally finite
discrepancies.

Fix a class \(C\), choose \(r\in C\), and partition it as \[
C=C^-\cup\{r\}\cup C^+,
\] where \[
C^-=\{x\in C:x\to r\},
\qquad
C^+=\{x\in C:r\to x\}.
\] For every \(x\in C\), the sets \(O(x)\) and \(O(r)\) differ finitely.
It follows that:

\begin{itemize}
\tightlist
\item
  every vertex of \(C^-\) has finite outdegree within \(C^-\);
\item
  every vertex of \(C^+\) has finite indegree within \(C^+\);
\item
  the pairs disagreeing with the orientation \(C^-\to C^+\) form a
  locally finite graph.
\end{itemize}

Enumerate \(C^-\) and \(C^+\), using finite enumerations when
appropriate. Order \(C^-\) in reverse enumeration order and \(C^+\) in
enumeration order, and concatenate them as \[
C^- \;<\; r \;<\; C^+.
\qquad\text{(3)}
\]

To check local finiteness within \(C^-\), consider its \(i\)-th
enumerated vertex. Among vertices enumerated before it, there are only
\(i\) possible discrepancies; among vertices enumerated after it,
discrepancies are contained in its finite outneighbourhood in \(C^-\).
The analogous argument for \(C^+\) uses finite indegree. Together with
the third bullet above and the exact orientations through \(r\), this
proves that the discrepancies inside \(C\) are locally finite.

Finally, order all vertices first by the chosen class order and then by
the orders (3). Let \(L\) be the resulting transitive tournament. By
(2), its cross-class discrepancies with \(S\) lie in the locally finite
graph \(J(S)\); its within-class discrepancies are locally finite by
construction. Each vertex therefore belongs to only finitely many
changed pairs. \(\square\)

\subsection{\texorpdfstring{3. Avoiding \(T_\ast\) forces \(J(S)\) to
have maximum degree at most
two}{3. Avoiding T\_\textbackslash ast forces J(S) to have maximum degree at most two}}\label{avoiding-t_ast-forces-js-to-have-maximum-degree-at-most-two}

For an arc \(x\to y\), also define \[
P(x,y)=\{p:x\to p\to y\}.
\] These are the vertices that would complete a directed triangle if
\(xy\) were reversed.

\subsubsection{Lemma 2}\label{lemma-2}

If \(S\) is \(T_\ast\)-free and \(x\to y\) has \(P(x,y)\ne\varnothing\),
then \[
|Z(x,y)|\le 2.
\]

\paragraph{Proof}\label{proof-1}

Choose \(p\in P(x,y)\). Suppose that \(Z(x,y)\) contains three vertices.
Among them, two, say \(z_1,z_2\), have the same orientation relative to
\(p\).

If \(p\to z_1,z_2\), then \[
\{y,p\}\ \longrightarrow\ \{z_1,z_2\}\
\longrightarrow\ \{x\}\ \longrightarrow\ \{y,p\}.
\] These five vertices induce \(T_\ast\).

If \(z_1,z_2\to p\), then \[
\{z_1,z_2\}\ \longrightarrow\ \{x,p\}\
\longrightarrow\ \{y\}\ \longrightarrow\ \{z_1,z_2\},
\] again inducing \(T_\ast\).

In either case, the two-vertex blocks automatically induce transitive
tournaments. Both possibilities contradict \(T_\ast\)-freeness.
\(\square\)

\subsubsection{Lemma 3}\label{lemma-3}

In any tournament, the arcs \(x\to y\) satisfying \[
P(x,y)=\varnothing
\qquad\text{(4)}
\] form an undirected graph of maximum degree at most two.

\paragraph{Proof}\label{proof-2}

At any vertex \(x\), at most one outgoing arc can satisfy (4). Indeed,
if both \(x\to y\) and \(x\to z\) satisfied it, then one of \(y\to z\)
and \(z\to y\) holds. In the first case, \[
x\to y\to z
\] contradicts \(P(x,z)=\varnothing\); the second case similarly
contradicts \(P(x,y)=\varnothing\).

Likewise, at most one incoming arc at \(x\) can satisfy (4). Thus the
corresponding undirected degree is at most two. \(\square\)

\subsubsection{Corollary 4}\label{corollary-4}

For every \(T_\ast\)-free tournament \(S\), \[
\Delta(J(S))\le 2.
\] Consequently, every countable \(T_\ast\)-free tournament is a locally
finite perturbation of a transitive tournament.

\paragraph{Proof}\label{proof-3}

If an arc belongs to infinitely many directed triangles, then \(Z(x,y)\)
is infinite. Lemma 2 forces \(P(x,y)=\varnothing\). Therefore \(J(S)\)
is a subgraph of the graph in Lemma 3, proving the degree bound. Lemma 1
gives the final assertion. \(\square\)

\subsection{4. Completing the disproof}\label{completing-the-disproof}

Let \(S\) be any countably infinite \(T_\ast\)-free tournament.

\textbf{Case 1: \(S\) is nontransitive.} By Corollary 4, there is a
transitive tournament \(L\) on \(V(S)\) whose change graph with \(S\) is
locally finite. Since \(S\) is nontransitive, \(L\ne S\). Since
\(T_\ast\) is nontransitive, \(L\) is \(T_\ast\)-free. Take \(S'=L\).

\textbf{Case 2: \(S\) is transitive.} Reverse any one arc to obtain
\(S'\). Every directed triangle in \(S'\) must contain that reversed
arc.

On the other hand, \(T_\ast\) contains the two edge-disjoint directed
triangles \[
c\to a_0\to b_0\to c
\qquad\text{and}\qquad
c\to a_1\to b_1\to c.
\] They share a vertex but no arc. Thus \(T_\ast\) cannot occur in a
tournament in which all directed triangles contain one fixed arc.
Therefore \(S'\) is \(T_\ast\)-free, and its change graph with \(S\)
consists of one edge.

Both cases produce a distinct \(T_\ast\)-free locally finite
perturbation of \(S\). Hence no tournament \(S_{T_\ast}\) as required by
the conjecture exists. \(\square\)

The counterexample addresses the full locally finite perturbation
formulation, including its intended nonempty-perturbation convention.
There is no remaining gap in this disproof; no assertion is made here
about the weaker requirement involving only individual arc reversals.

\end{document}
